% -*- coding: utf-8 -*-
% =============================================================================
% header.tex —— 竞赛少题深研学案（v2 单源多视图版）：导言区 + 宏定义 + 标题区 + 双栏开始
% 架构: <入口>.tex -> main.tex -> [header.tex, sec*.tex, footer.tex]
%       （sec1.tex 等切片即内容块装配写作层，按 materials.json 块属性装配）
% 入口: student.tex（student_initial 视图）/ teacher.tex（teacher 视图）/
%       hint-<pid>.tex / review-<pid>.tex（gen_config.py 按 config.yaml views 生成，
%       无对应内容块时不生成空入口）。
% 参数: 全部版面与文案参数来自 config-class.tex / config.tex
%       （由 scripts/gen_config.py 依据 config.yaml 自动生成，勿手改）。
% =============================================================================

% ---- 文档类选项（字号等，须在 \documentclass 之前注入） ----
\input{config-class.tex}
\documentclass{ctexart}

% ---- 版面几何与文本参数 + 视图开关（config.yaml 自动生成） ----
% config.tex 依据入口文件声明的 \pdCurrentView 定义 \pdView@<视图ID> 开关；
% 未声明 \pdCurrentView 的旧入口（兼容模式）不定义任何视图开关。
\input{config.tex}

% ---- 当前视图名（写入日志 SUMMARY，供 build.py 逐视图核对） ----
\ifdefined\pdCurrentView
  \providecommand{\pdCurrentViewName}{}%
  \renewcommand{\pdCurrentViewName}{\pdCurrentView}%
\else
  \ifdefined\TeacherVersion
    \newcommand{\pdCurrentViewName}{teacher}%
  \else
    \newcommand{\pdCurrentViewName}{student}%
  \fi
\fi

% ---- 教师版版面覆盖：A4 竖版、单栏 ----
% 学生版沿用 config.yaml 的大8开横版双栏；教师版供案头翻阅，固定为 A4 竖版单栏。
% \geometry 在导言区可重复调用并覆盖先前设置；\pdColumns=1 时 \pdBeginColumns 不产生 multicols 环境。
\ifdefined\TeacherVersion
  \geometry{a4paper, portrait, top=2.5cm, bottom=2.5cm, left=2.5cm, right=2.5cm, headheight=25pt}
  \renewcommand{\pdColumns}{1}
\fi

% ---- 学生解答版版面覆盖：A4 竖版、单栏 ----
\ifdefined\StudentSolutionVersion
  \geometry{a4paper, portrait, top=2.5cm, bottom=2.5cm, left=2.5cm, right=2.5cm, headheight=25pt}
  \renewcommand{\pdColumns}{1}
\fi

% ---- 独立提示/复盘页（separate_hint / separate_review 渠道入口）：A4 竖版单栏 ----
% 由生成的 hint-*/review-* 入口文件 \def\pdCompactPage{} 触发。
\ifdefined\pdCompactPage
  \geometry{a4paper, portrait, top=2.5cm, bottom=2.5cm, left=2.5cm, right=2.5cm, headheight=25pt}
  \renewcommand{\pdColumns}{1}
\fi

% ---- 栏环境起止（统一由此处负责，main.tex/footer.tex 只调用宏） ----
% multicols 不会因参数为 1 自动退化为单栏（会警告并强制改为双栏），
% 因此栏数为 1 时不得进入 multicols 环境。
\newcommand{\pdBeginColumns}{%
  \ifnum\pdColumns>1 \begin{multicols}{\pdColumns}\fi
}
\newcommand{\pdEndColumns}{%
  \ifnum\pdColumns>1 \end{multicols}\fi
}

% -----------------------------------------------------------------------------
% 核心宏包加载
% -----------------------------------------------------------------------------
\usepackage{amsmath,amssymb,amsthm,bm,mathtools}
\usepackage{fancyhdr}
\usepackage{multicol}
\usepackage{enumitem}
\usepackage{graphicx}
\usepackage{xcolor}
\usepackage{comment}
\usepackage[framemethod=TikZ]{mdframed}
\usepackage{tikz}
\usetikzlibrary{arrows.meta, positioning, calc, shapes.geometric}
\usepackage{exam-zh-choices}  % 选择题选项自动排版宏包
\usepackage{hyperref}
\usepackage{lastpage}

\hypersetup{
  colorlinks=true,
  linkcolor=blue!75!black,
  citecolor=blue!75!black,
  urlcolor=blue!75!black
}

% -----------------------------------------------------------------------------
% 统计计数器（编译后输出环境使用摘要，便于验收与排错）
% -----------------------------------------------------------------------------
\newcounter{pdSolutionCount}
\newcounter{pdNoteCount}
\newcounter{pdFillinCount}
\newlength{\fillinwidth}

% -----------------------------------------------------------------------------
% 视图条件装配宏（v2 单源多视图核心机制）
% 当前编译入口的视图由入口文件 \def\pdCurrentView{<id>} 声明，config.tex 据此
% 定义 \pdView@<id> 开关（未在 config.yaml views 声明的视图直接编译报错）。
%
%   \ViewBlock{<views 列表>}{<内容>}
%       内容仅在当前视图属于 views 列表时输出，否则整段丢弃。
%       materials 模式（teaching/materials.json 存在）下，切片中每个内容块
%       以 `% block: <block_id>` 注释开头，其后 \ViewBlock 的 views 列表必须
%       与 materials.json 中该块的 views 完全一致（check_numbering.py 核对）。
%
% 机制性防泄露（对应 diagnoses.md 揭示层级失效分析）：
%   - 未声明 \pdCurrentView 的编译不定义任何视图开关，\ViewBlock 内容一律
%     不输出——视图未声明时宁可不印，也不误印；
%   - knowledgebox / hintbox / reviewbox 等含结论性内容的环境默认仅 teacher
%     视图可见，学生视图必须显式声明（见下方各环境定义）；
%   - available_after 不是 initial 的块在静态 PDF 中不会因为"印在留白之后"
%     而延迟揭示，只能经 separate_hint / separate_review 独立入口分阶段发放。
% -----------------------------------------------------------------------------
\makeatletter
\def\pd@yes{1}
% \pdIfViewIn{<views 列表>}{<命中时内容>}{<未命中时内容>}
\newcommand{\pdIfViewIn}[3]{%
  \begingroup
    \edef\pd@list{\zap@space#1 \@empty}%
    \let\pd@hit\relax
    \@for\pd@v:=\pd@list\do{%
      \expandafter\ifx\csname pdView@\pd@v\endcsname\relax\else
        \let\pd@hit\pd@yes
      \fi}%
    \ifx\pd@hit\pd@yes
      \def\pd@next{\endgroup#2}%
    \else
      \def\pd@next{\endgroup#3}%
    \fi
  \pd@next}
\makeatother

% 视图条件内容块：\ViewBlock{student_initial,teacher}{内容}
\newcommand{\ViewBlock}[2]{\pdIfViewIn{#1}{#2}{}}

% -----------------------------------------------------------------------------
% 双版本环境控制宏 (Single-source Dual-output)
% 说明: solution/teacherNote/fillin/ansspace 在两版中的行为在此处统一定义，
%       切片文件 (sec*.tex) 只需使用同一套命令，无需关心版本差异。
%       教师版判定沿用 \TeacherVersion 开关（teacher 视图入口声明）；
%       hint/review 入口不声明 \TeacherVersion，按学生侧规则剥离
%       solution/teacherNote（双保险：内容块本应由 \ViewBlock 按视图装配）。
% -----------------------------------------------------------------------------
\ifdefined\TeacherVersion
  % ============================ 教师版环境配置 ============================
  \pagestyle{fancy}
  \fancyhf{}
  \fancyhead[L]{\bfseries\pdTeacherHeadLeft}
  \fancyfoot[C]{\thepage\ of \pageref{LastPage}}
  \renewcommand{\headrulewidth}{0.4pt}
  \renewcommand{\footrulewidth}{0.4pt}

  % 题型标签：题号后不带【】说明文字，保持版面干净
  \newcommand{\qtype}[1]{}

  % 双栏折页排版宏（教师单栏版不强制分栏或分页）
  \newcommand{\pdColumnBreak}{}
  \newcommand{\pdPageBreak}{}

  % 题目状态标签 \practicemode{guided}（可选，缺省 guided）：
  % 教学模式唯一取值为 guided（分步引导）；取值非法时发出编译警告。
  \providecommand{\practicemode}[1]{%
    \def\pdPmode{#1}%
    \def\pdPmodeG{guided}%
    \ifx\pdPmode\pdPmodeG\else
      \PackageWarning{problem-driven}{\protect\practicemode\space 取值只能是 guided（收到“#1”）}%
    \fi}

  % 参考解答环境（显示完整答案推导，并计数）
  \newenvironment{solution}[1][解]{%
    \stepcounter{pdSolutionCount}%
    \par\begingroup%
    \begin{proof}[\indent\bfseries【#1】]%
  }{%
    \end{proof}\endgroup\par
  }

  % 教师专用教学要点、易错警示与评分提示环境
  \newenvironment{teacherNote}[1][教学要点与易错分析]{%
    \stepcounter{pdNoteCount}%
    \par\vspace{0.6ex}%
    \begingroup\small
    \noindent\textbf{【#1】}\par\vspace{0.2ex}%
  }{%
    \endgroup\par\vspace{0.6ex}%
  }

  % 答题留白：教师版压缩为微小行距，消除多余空白
  \newcommand{\ansspace}[1]{\par\vspace{0.4ex}}

  % 填空宏：教师版在下划线居中呈现醒目的加粗答案（黑白印刷友好）；
  % 缺宽度或缺答案时发出编译警告
  \newcommand{\fillin}[2][]{%
    \stepcounter{pdFillinCount}%
    \if\relax\detokenize{#2}\relax
      \PackageWarning{problem-driven}{\protect\fillin\space 缺少预留宽度参数（第 \thepdFillinCount\space 处）}%
    \fi
    \if\relax\detokenize{#1}\relax
      \PackageWarning{problem-driven}{\protect\fillin\space 未提供参考答案（第 \thepdFillinCount\space 处）}%
      \underline{\hspace{#2}}%
    \else
      \ifmmode
        \settowidth{\fillinwidth}{$\bm{#1}$}%
        \ifdim\fillinwidth>#2
          \underline{\hspace{0.5em}{\bm{#1}}\hspace{0.5em}}%
        \else
          \underline{\makebox[#2][c]{$\bm{#1}$}}%
        \fi
      \else
        \settowidth{\fillinwidth}{\bfseries#1}%
        \ifdim\fillinwidth>#2
          \underline{\hspace{0.5em}\textbf{#1}\hspace{0.5em}}%
        \else
          \underline{\makebox[#2][c]{\bfseries#1}}%
        \fi
      \fi
    \fi
  }
\else
  \ifdefined\StudentSolutionVersion
    % ==================== 学生参考解答版环境配置 ====================
    \pagestyle{fancy}
    \fancyhf{}
    \fancyhead[L]{\bfseries\ifdefined\pdMainTitle\pdMainTitle\else 学案\fi\quad 参考解答}
    \fancyfoot[C]{\thepage\ of \pageref{LastPage}}
    \renewcommand{\headrulewidth}{0.4pt}
    \renewcommand{\footrulewidth}{0.4pt}

    % 排版优化：疏松自然的段落间距与行距，公式间距舒适
    \setlength{\parskip}{0.8ex plus 0.3ex minus 0.2ex}
    \linespread{1.18}\selectfont
    \setlength{\abovedisplayskip}{7pt plus 2pt minus 2pt}
    \setlength{\belowdisplayskip}{7pt plus 2pt minus 2pt}

    \newcommand{\qtype}[1]{}
    \newcommand{\pdColumnBreak}{\par\vspace{1.5ex}}
    \newcommand{\pdPageBreak}{\clearpage}
    \providecommand{\practicemode}[1]{}

    % 参考解答环境：显示解答，段前段后适度留白，标题显式突出
    \newenvironment{solution}[1][参考解答]{%
      \stepcounter{pdSolutionCount}%
      \par\vspace{0.8ex}\begingroup
      \begin{proof}[\indent\bfseries【#1】]%
    }{%
      \end{proof}\endgroup\vspace{0.8ex}\par
    }
    % 教师批注仍剥离
    \specialcomment{teacherNote}{\stepcounter{pdNoteCount}\begingroup\def\ProcessCutFile{}}{\endgroup}
    % 答题留白在解答版中忽略
    \newcommand{\ansspace}[1]{}
    \newcommand{\fillin}[2][]{%
      \stepcounter{pdFillinCount}%
      \if\relax\detokenize{#1}\relax
        \underline{\hspace{#2}}%
      \else
        \textbf{#1}%
      \fi
    }
  \else
    % ==================== 学生侧环境配置（含 hint/review 入口） ====================
    \pagestyle{fancy}
    \fancyhf{}
    \fancyhead[L]{姓名：\underline{\hspace{2.2cm}}\quad 班级：\underline{\hspace{2cm}}}
    \fancyfoot[C]{\thepage\ of \pageref{LastPage}}
    \renewcommand{\headrulewidth}{0.4pt}
    \renewcommand{\footrulewidth}{0.4pt}

    % 题型标签：学生版隐藏【xx题】字样（设计流程保留题型，但不明示）
    \newcommand{\qtype}[1]{}

    % 双栏折页排版宏（学生版梳松排版：按栏换栏、按页换页）
    \newcommand{\pdColumnBreak}{\columnbreak}
    \newcommand{\pdPageBreak}{\newpage}

    % 题目状态标签：学生版隐藏（引导隐藏由切片写作纪律保证，
    % 并由 scripts/check_numbering.py 检查题块内不得含提示标记）
    \providecommand{\practicemode}[1]{}%

    % 彻底剥离答案与教师批注，同时计数以便编译日志中核对剥离数量
    % 注意: \specialcomment 默认会“包含”环境体；在组内将 \ProcessCutFile 重定义为
    % 空即恢复剥离效果（comment 宏包官方推荐用法），begin 行尾内容会被自动丢弃。
    \specialcomment{solution}{\stepcounter{pdSolutionCount}\begingroup\def\ProcessCutFile{}}{\endgroup}
    \specialcomment{teacherNote}{\stepcounter{pdNoteCount}\begingroup\def\ProcessCutFile{}}{\endgroup}

    % 答题留白：学生版产生实际书写所需的物理垂直空间
    \newcommand{\ansspace}[1]{\vspace{#1}}

    % 填空宏：学生版仅呈现预留长度的下划线空白
    \newcommand{\fillin}[2][]{%
      \stepcounter{pdFillinCount}%
      \if\relax\detokenize{#2}\relax
        \PackageWarning{problem-driven}{\protect\fillin\space 缺少预留宽度参数（第 \thepdFillinCount\space 处）}%
      \fi
      \underline{\hspace{#2}}%
    }
  \fi
\fi

% 编译结束时的质量摘要（写入 .log 与终端，供 build.py 解析核对）
% SUMMARY 行保持旧版格式（不超过 79 列防 \typeout 折行）；视图名单独一行 VIEW。
\AtEndDocument{%
  \typeout{[problem-driven] VIEW \pdCurrentViewName}%
  \typeout{[problem-driven] SUMMARY solution=\arabic{pdSolutionCount},
    teacherNote=\arabic{pdNoteCount}, fillin=\arabic{pdFillinCount}}%
  \ifdefined\TeacherVersion
    \ifnum\value{pdSolutionCount}=0
      \PackageWarning{problem-driven}{教师版未发现任何 solution 环境，请检查切片文件}%
    \fi
  \else
    \ifdefined\StudentSolutionVersion
      \typeout{[problem-driven] 学生解答版包含 solution=\arabic{pdSolutionCount}}%
    \else
      \typeout{[problem-driven] 学生版已剥离 solution/teacherNote（视图 \pdCurrentViewName），数量为上述 SUMMARY}%
    \fi
  \fi
}

% -----------------------------------------------------------------------------
% 微型知识点模块与列表规范
% \begin{microknowledge}[标题][可见视图列表]
%   省略视图列表时按 L0 前置知识处理，全部视图可见；
%   显式给出视图列表（如 [教师版预备知识][teacher]）时仅所列视图可见。
%   注意：microknowledge 的内容层级由 materials 的块属性与 check_teaching /
%   check_student_visibility / check_student_task 检查约束，写作时不得把 L1/L2 策略与结论包装成"知识卡"。
% -----------------------------------------------------------------------------
\makeatletter
\newsavebox{\microbox}
\newcommand{\pd@microrender}[2]{%
  \par\vspace{0.8ex}%
  \begin{lrbox}{\microbox}%
    \begin{minipage}{\dimexpr\linewidth-2\fboxsep-2\fboxrule\relax}%
      \textbf{【#1】}\par\vspace{0.3ex}\small #2%
    \end{minipage}%
  \end{lrbox}%
  \noindent\fbox{\usebox{\microbox}}\par\vspace{0.8ex}%
}
\NewDocumentEnvironment{microknowledge}{O{核心定义与记号} o +b}{%
  \IfNoValueTF{#2}%
    {\pd@microrender{#1}{#3}}%
    {\edef\pd@views{\zap@space#2 \@empty}%
     \pdIfViewIn{\pd@views}{\pd@microrender{#1}{#3}}{}}%
}{}
\makeatother

% -----------------------------------------------------------------------------
% 结论框环境（虚线框）：默认仅 teacher 视图可见（机制性防泄露）
% \begin{knowledgebox}[标题][追加可见视图]
%   省略第二可选参数时仅 teacher 视图可见；学生侧视图（student_initial /
%   hint-* / review-*）必须显式声明，如 [本节小结][review-p1]。
%   v2 不再要求每切片末必须有结论框；结论框只放经审定、允许在该视图出现的
%   内容（diagnoses.md：旧版 knowledgebox 无版本隔离导致学生直接看到答案）。
% \begin{teacherKnowledgebox}[标题]  —— 固定仅 teacher 视图的结论/复盘框。
% \begin{hintbox}[标题][追加可见视图] —— 提示框，默认仅 teacher，提示页须
%   显式声明对应 hint-* 视图（或由装配层 \ViewBlock 统一控制）。
% \begin{reviewbox}[标题][追加可见视图] —— 复盘框，默认仅 teacher，复盘页
%   须显式声明对应 review-* 视图。
% -----------------------------------------------------------------------------
\mdfdefinestyle{pdKnowledgeBoxStyle}{%
  linecolor=black,
  linewidth=0.6pt,
  backgroundcolor=white,
  draw=black,  % 移除彩色设置
  dash pattern=on 3pt off 2pt,
  roundcorner=0pt,
  innertopmargin=6pt,
  innerbottommargin=6pt,
  innerleftmargin=8pt,
  innerrightmargin=8pt,
  nobreak=true
}
\mdfdefinestyle{pdHintBoxStyle}{%
  linecolor=black,
  linewidth=0.6pt,
  backgroundcolor=white,
  draw=black,  % 移除彩色设置
  dash pattern=on 2pt off 2pt,
  roundcorner=0pt,
  innertopmargin=6pt,
  innerbottommargin=6pt,
  innerleftmargin=8pt,
  innerrightmargin=8pt,
  nobreak=true
}
\mdfdefinestyle{pdReviewBoxStyle}{%
  linecolor=black,
  linewidth=0.6pt,
  backgroundcolor=white,
  roundcorner=0pt,
  innertopmargin=6pt,
  innerbottommargin=6pt,
  innerleftmargin=8pt,
  innerrightmargin=8pt,
  nobreak=true
}
\makeatletter
% #1=样式 #2=标题 #3=正文
\newcommand{\pd@boxrender}[3]{%
  \par\vspace{0.8ex}%
  \begin{mdframed}[style=#1]%
  \textbf{#2}\par\vspace{0.3ex}\small #3%
  \end{mdframed}\par\vspace{0.8ex}%
}
% 默认视图列表 = teacher；显式第二可选参数为追加视图
\NewDocumentEnvironment{knowledgebox}{O{本节结论} o +b}{%
  \IfNoValueTF{#2}%
    {\def\pd@views{teacher}}%
    {\edef\pd@views{teacher,\zap@space#2 \@empty}}%
  \pdIfViewIn{\pd@views}{\pd@boxrender{pdKnowledgeBoxStyle}{#1}{#3}}{}%
}{}
\NewDocumentEnvironment{teacherKnowledgebox}{O{教师结论与复盘} +b}{%
  \pdIfViewIn{teacher}{\pd@boxrender{pdKnowledgeBoxStyle}{#1}{#2}}{}%
}{}
\NewDocumentEnvironment{hintbox}{O{提示} o +b}{%
  \IfNoValueTF{#2}%
    {\def\pd@views{teacher}}%
    {\edef\pd@views{teacher,\zap@space#2 \@empty}}%
  \pdIfViewIn{\pd@views}{\pd@boxrender{pdHintBoxStyle}{#1}{#3}}{}%
}{}
\NewDocumentEnvironment{reviewbox}{O{复盘} o +b}{%
  \IfNoValueTF{#2}%
    {\def\pd@views{teacher}}%
    {\edef\pd@views{teacher,\zap@space#2 \@empty}}%
  \pdIfViewIn{\pd@views}{\pd@boxrender{pdReviewBoxStyle}{#1}{#3}}{}%
}{}
\makeatother


% 列表样式统一设定 (一级大题号连续递增 1. 2. 3.；二级小问题号用 (1), (2), (3)...)
\setlist[enumerate,1]{label=\textbf{\arabic*.}, leftmargin=*, itemsep=0.8em, topsep=0.4ex}
\setlist[enumerate,2]{label=(\arabic*), leftmargin=*, itemsep=0.4em}

% -----------------------------------------------------------------------------
% 题链自动编号环境 probchain（推荐）
% 自建计数器跨环境续号，彻底摆脱手工维护 series=/resume= 的对位负担：
% 所有小节统一写作 \begin{probchain}...\end{probchain}，先后顺序天然决定题号。
% 旧写法 \begin{enumerate}[series=probchain] / [resume=probchain] 仍然兼容。
% 注意：probchain 主体应按视图整体装配（如 \ViewBlock{student_initial,teacher}），
% 避免提示/复盘视图渲染出没有任何 \item 的空题链（空 enumerate 是编译错误）。
% -----------------------------------------------------------------------------
\newcounter{pdprob}
\newenvironment{probchain}{%
  \begin{enumerate}[start=\the\numexpr\value{pdprob}+1\relax]%
}{%
    \setcounter{pdprob}{\value{enumi}}%
  \end{enumerate}%
}
